\section{Actual positive member histories under smooth pullback}
\label{hir:smooth-positive-member-history-section}
This extension supplies the positive-boundary branch needed in the smooth
comparison of the original recursive invariant. It uses the literal marked
ideal, the original indexed divisor, the actual coefficient-homogenized ideal
$\overline J$, and actual Cartesian test histories. The zero-boundary
codimension-one branch, chosen tangent-direction masks, and the complete
higher recursive comparison remain separate obligations. No equality of
chosen tangent directions or of the entire higher invariant is assumed. This is the smooth-map extension of the
original Step One construction in Wlodarczyk,
proof of the marked-ideal resolution proposition, Step 1aa/1ab.

\begin{lemma}[The actual coefficient-homogenized record]
\label{hir:smooth-bar-record}
\source{wlodarczyk05-hironaka,bierstone-milman08-functoriality}
\lean{Hironaka.HironakaSmoothSubmersionPositiveMemberHistory.pullback_barMarked,Hironaka.HironakaSmoothSubmersionPositiveMemberHistory.IsActualTowerPullback.bar_stage_dictionary,Hironaka.HironakaSmoothSubmersionPositiveMemberHistory.IsActualTowerPullback.bar_original_members_dictionary,Hironaka.HironakaSmoothSubmersionPositiveMemberHistory.IsActualTowerPullback.bar_level_dictionary}
\uses{hir:submersion-relative-differential-whole-homogenized-coefficient,hir:submersion-relative-control-joint-tower,hir:submersion-cartesian-tower-actual-history}
\leanok
For a smooth holomorphic map $f:Y\to X$ and an original well-formed marked
ideal $J$, $f^*\overline J=\overline{f^*J}$. If $Q\to T$ is an actual
Cartesian history and $T$ is a geometric test history for $\overline J$,
then $Q$ is a test history for $\overline{f^*J}$ and their entire transformed
records correspond under the actual last-stage smooth map $g$.
\end{lemma}
\begin{proof}
Relative derivative ideals, homogenization and coefficient ideals commute
with smooth pullback, including the actual carrier, marking and full divisor
list. Apply those two identities successively to $\overline J=C_N(H_N(J))$.
The established geometric Cartesian controlled-transform identity, applied
to the well-formed $\overline J$, gives both the test-history assertion and
the transformed record. The original $|E|$ strict transforms are the initial
$|E|$ entries of that transformed record, so their lists and boundary counts
are exactly pulled back.
\end{proof}

\begin{lemma}[Actual member locus and positive special branch]
\label{hir:smooth-positive-member-special}
\source{wlodarczyk05-hironaka,bierstone-milman08-functoriality}
\lean{Hironaka.HironakaSmoothSubmersionPositiveMemberHistory.IsActualTowerPullback.bar_member_locus_dictionary,Hironaka.HironakaSmoothSubmersionPositiveMemberHistory.IsActualTowerPullback.positive_isSpecial_dictionary,Hironaka.HironakaSmoothSubmersionPositiveMemberHistory.IsActualTowerPullback.positive_stepRegime_dictionary,Hironaka.HironakaSmoothSubmersionPositiveMemberHistory.IsActualTowerPullback.positive_stepRegime_run}
\uses{hir:smooth-bar-record}
\leanok
At every stage, the original member-intersection locus of the cosupport of
$\overline{f^*J}$ is the inverse image of the literal original member locus
for $\overline J$. If the target boundary count at $g(y)$ is positive,
$y$ is special for the source Step One history if and only if $g(y)$ is
special for the target history. For every fixed positive boundary count
$s$, the actual Step One regime at level $s$ corresponds at all points,
and its complete recursive run index is unchanged.
\end{lemma}
\begin{proof}
A smooth map is continuous and open; thus it maps the neighbourhood filter
at $y$ onto the neighbourhood filter at $g(y)$. In the definition of the
member locus, transport the cosupport, carrier, full original divisor list,
and the quantified nearby member-intersection condition with this exact
filter equality. The special branch at positive boundary count is precisely
membership in this member locus; the level-zero alternative is false.
For a fixed positive $s$, points whose level differs from $s$ satisfy neither
regime, and the remaining points use the preceding special-point comparison.
Induction on the actual history compares the regime at every intermediate
stage and therefore the recursive run index.
\end{proof}

\begin{lemma}[The actual positive member mask]
\label{hir:smooth-positive-member-mask}
\source{wlodarczyk05-hironaka,bierstone-milman08-functoriality}
\lean{Hironaka.HironakaSmoothSubmersionPositiveMemberHistory.IsActualTowerPullback.positive_maskA_dictionary}
\uses{hir:smooth-positive-member-special}
\leanok
For any point of positive boundary count, the literal Step One member mask
$A_{f^*J,Q}(y)$ is $g^{-1}A_{J,T}(g(y))$.
\end{lemma}
\begin{proof}
The condition excluding every original member missing the base point
corresponds by the original list dictionary. In the remaining condition,
if a point has the same boundary count as the base point, this count is
positive, and the preceding special-point comparison applies; if the counts
differ, both exclusions are false. The raw sets in the two masks therefore
correspond exactly. For a continuous open map, inverse image commutes with
interior, proving the statement for the actual open masks. No chosen contact
mask is substituted for these masks.
\end{proof}
